%!TEX root = ../main.tex
\thispagestyle{empty}

\begin{center}
\Large
\textbf{誌\hspace{1cm}謝}
\end{center}

\vspace{1cm}
\linespread{2}%
\selectfont
\hspace{0.25cm}

% 下面開始寫致謝內容
本論文能夠順利完成，首先要感謝我的指導教授，在研究方向、問題定義、方法設計與論文撰寫上給予我許多寶貴的建議與指導。從最初對 RF 調變分類與對抗式攻擊防禦問題的探索，到後續方法設計、實驗分析與結果整理，老師的指導讓我能夠更清楚地理解研究問題，也使本論文更加完整。

同時，我也要感謝口試委員們在論文審查過程中所提供的寶貴意見。各位委員的建議使我能夠從不同角度重新檢視研究方法、實驗設計與論文論述，並進一步改善本研究的完整性與嚴謹性。

感謝實驗室的同學們，在研究過程中與我討論技術問題、分享實驗經驗，並在我遇到困難時給予協助與鼓勵。無論是 RF signal processing、deep learning model、adversarial attack，或是實驗環境建置上的討論，都對本論文的完成有很大的幫助。

最後，我要感謝我的家人與朋友，在這段研究期間給予我支持、包容與鼓勵。研究過程中難免遇到挫折與壓力，但因為有你們的陪伴與支持，讓我能夠持續前進並完成這篇論文。

謹以此論文獻給所有曾經幫助、支持與鼓勵我的人。

\vspace{3cm}
\begin{flushright}
\studentCnName \ 於

國立陽明交通大學\ \NameofDepartmentInstituteCN

\ThesisDateTW
\end{flushright}
